\section{RLHF 数据收集}

\subsection{成对反馈的标准流程}

\begin{figure}[H]
\centering
\includegraphics[width=0.95\textwidth]{slides-images/slide-035.jpg}
\caption{成对反馈收集的标准界面：标注者看到两个 AI 回复，选择哪个更好\protect\footnotemark}
\end{figure}
\footnotetext{Slides 第36页。}

成对反馈的基本流程是：
\begin{enumerate}
    \item 给定一个提示（prompt），模型生成多个不同的回复
    \item 标注者比较这些回复，标记哪个更好
    \item 收集大量这样的成对比较数据
\end{enumerate}

\subsection{标注指南的设计}

\begin{figure}[H]
\centering
\includegraphics[width=0.95\textwidth]{slides-images/slide-036.jpg}
\caption{InstructGPT 的标注指南：评估维度包括有用性（Helpful）、真实性（Truthful）和无害性（Harmless）\protect\footnotemark}
\end{figure}
\footnotetext{Slides 第37页。来自 InstructGPT 论文附录。}

InstructGPT 的标注指南定义了三个评估维度（HHH）：

\begin{itemize}
    \item \textbf{有用性（Helpful）}：使用清晰语言、回答用户真正想问的问题、对国际化敏感（如``football''不应默认为美式足球）
    \item \textbf{真实性（Truthful）}：不产生幻觉输出
    \item \textbf{无害性（Harmless）}：不包含有毒内容、不生成 NSFW 内容
\end{itemize}

\begin{figure}[H]
\centering
\includegraphics[width=0.95\textwidth]{slides-images/slide-037.jpg}
\caption{泄露的 Google Bard 标注指南：类似的有用性、风格、安全性评估维度，标注者每个问题仅有约一分钟\protect\footnotemark}
\end{figure}
\footnotetext{Slides 第38页。}

InstructGPT 通过 Scale 和 Upwork 平台收集了约 40 名标注者的数据。Google Bard 的指南包含类似的维度（有用性、风格、安全性），但据报道标注者每个问题仅有约一分钟的判断时间。

\subsection{众包标注的挑战}

\begin{figure}[H]
\centering
\includegraphics[width=0.95\textwidth]{slides-images/slide-040.jpg}
\caption{众包数据收集面临的三大挑战\protect\footnotemark}
\end{figure}
\footnotetext{Slides 第41页。}

课堂互动实验生动地展示了成对标注的困难。在约 5 分钟的时间内，斯坦福学生需要对多组回复进行比较判断：

\begin{warningbox}{成对标注中的三大陷阱}
\begin{enumerate}
    \item \textbf{难以获得高质量、可验证的标注者}：即便是有动力的学生也会出错
    \item \textbf{难以让标注者真正核查事实正确性}：大多数学生在 5 分钟内无法逐一验证回复中的所有事实声称
    \item \textbf{GPT-4 作弊问题}：一些标注者会直接将任务丢给 GPT-4，导致标注数据反映的是 AI 的偏好而非人类的真实判断
\end{enumerate}
\end{warningbox}

实验的一个关键发现是：许多更长的回复实际上包含了幻觉内容，但大多数人仍然选择了更长的版本。这直接说明了为什么长度偏差在 RLHF 数据中如此普遍。

\subsection{标注者偏见与多样性问题}

\begin{figure}[H]
\centering
\includegraphics[width=0.95\textwidth]{slides-images/slide-042.jpg}
\caption{标注者背景对模型价值观的影响：InstructGPT 的标注者主要来自菲律宾和孟加拉国\protect\footnotemark}
\end{figure}
\footnotetext{Slides 第43页。}

RLHF 处于训练流程的末端，因此对模型最终行为有极强的影响力。Tatsu 团队的研究发现，InstructGPT 在训练后变得更加认同东南亚宗教的观点——而查看论文附录后发现，标注者主要来自菲律宾和孟加拉国（仅 17\% 为美国人）。这一发现提醒我们：

\begin{warningbox}{标注者构成决定模型价值观}
由于 RLHF 在流程末端、对模型行为影响最大，标注者的文化背景、教育水平、个人偏好会直接影响模型的价值观取向。不同标注者群体关注的维度也不同——研究者/作者更关注事实正确性，而众包工人更关注格式美观度。
\end{warningbox}

\begin{figure}[H]
\centering
\includegraphics[width=0.95\textwidth]{slides-images/slide-043.jpg}
\caption{不同标注者群体的关注点差异：研究者（Authors）更关注事实性，众包工人（Crowd workers）更关注格式\protect\footnotemark}
\end{figure}
\footnotetext{Slides 第44页。来自 Hosking, Blunam, \& Barlo 的研究。}

\subsection{AI 反馈：从 RLHF 到 RLAIF}

\begin{figure}[H]
\centering
\includegraphics[width=0.95\textwidth]{slides-images/slide-045.jpg}
\caption{AI 反馈的有效性：GPT-4 与人类评审的一致性与人类之间的一致性相当\protect\footnotemark}
\end{figure}
\footnotetext{Slides 第46页。}

鉴于人类标注的种种挑战，AI 反馈（RLAIF）成为了一个越来越流行的替代方案。核心发现包括：

\begin{itemize}
    \item GPT-4 的成对偏好判断与人类评审的一致性（Agreement），与人类之间的相互一致性大致相当
    \item AI 反馈的成本远低于人类标注
    \item Hugging Face 在构建 Zephyr 7B 模型时，最初坚持使用人类众包数据，但最终发现 GPT-4 生成的反馈效果更好
\end{itemize}

\begin{figure}[H]
\centering
\includegraphics[width=0.95\textwidth]{slides-images/slide-047.jpg}
\caption{Tulu 3 的数据构建流程：使用不同模型生成回复，再由 LM 进行成对偏好判断\protect\footnotemark}
\end{figure}
\footnotetext{Slides 第48页。来自 AI2 Tulu 3 项目。}

这一趋势可以追溯到 Anthropic 的 Constitutional AI 论文，该工作首次系统性地提出用 AI 反馈替代人类反馈进行对齐训练。现代实践中，如 Tulu 3（AI2）等项目，已经完全采用 LM 生成的反馈数据。

\begin{knowledgebox}{AI 反馈的注意事项}
\begin{itemize}
    \item \textbf{自我偏好偏差}：大多数模型对自己的输出有可检测的偏好，使用同一模型既生成又评判时需要特别小心
    \item \textbf{长度偏差}：AI 评审同样偏好更长的回复，可能导致模型输出越来越冗长
    \item \textbf{风格偏好}：AI 评审可能将特定风格（如列表格式）与质量混淆
\end{itemize}
\end{knowledgebox}

\subsection{本章小结}

\begin{importantbox}{RLHF 数据收集的核心教训}
\begin{enumerate}
    \item 成对反馈虽然比 SFT 数据更容易收集，但仍面临事实核查困难、长度偏差、标注者偏见等严峻挑战
    \item 标注者背景对模型价值观有直接影响，需要审慎考虑标注团队的多样性和代表性
    \item AI 反馈已经成为主流替代方案，但自身也带来新的偏差问题
\end{enumerate}
\end{importantbox}


